% =========================================================================
% PRÁCTICA 8: SECUENCIADOR MEDIANTE BARRA DE LEDS
% =========================================================================
\section{Práctica 8: Secuenciador mediante Barra de LEDs}

\subsection{Descripción General}
En este desarrollo experimental se diseñó e implementó un sistema de control de efectos lumínicos empleando el microcontrolador PIC16F877A acoplado a un arreglo lineal de diodos (Barra gráfica de LEDs). El propósito principal radica en la gestión de flujos de salida digitales mediante la monitorización activa (\textit{polling}) de terminales de entrada discretas. A través de la manipulación de registros y máscaras de bits, el sistema conmuta en tiempo real entre un estado de reposo y tres rutinas dinámicas de iluminación independientes.

\subsection{Componentes Utilizados}
El hardware seleccionado para la implementación física comprende los siguientes elementos:
\begin{itemize}
    \item \textbf{Elementos de Control y Procesamiento:} 1 Microcontrolador PIC16F877A, 1 Programador PICkit 3.
    \item \textbf{Frecuencia y Estabilización:} 1 Cristal de cuarzo de $4\text{ MHz}$, 2 Capacitores cerámicos de $22\text{ pF}$.
    \item \textbf{Interfaz de Entrada/Salida:} 1 Barra de LEDs, 1 Pulsador (\textit{Push button}).
    \item \textbf{Red de Resistencias:} 4 Resistencias de $220\ \Omega$, 1 Resistencia de $1\text{ k}\Omega$.
    \item \textbf{Soporte Físico:} 1 Placa de pruebas (\textit{Protoboard}) y conductores eléctricos variados.
\end{itemize}

\subsection{Entorno de Desarrollo Software}
\begin{itemize}
    \item \textbf{IDE:} MPLAB X IDE v6.25.
    \item \textbf{Compilador:} MPLAB XC8 Compiler.
    \item \textbf{Entorno de Simulación:} Proteus 8.
\end{itemize}

\subsection{Análisis Operativo y Arquitectura Lógica}
El algoritmo base rige su comportamiento de acuerdo con el estado lógico de los pines menos significativos del Puerto A (\texttt{RA0}, \texttt{RA1} y \texttt{RA2}). Al deshabilitar los módulos de comparación analógica a través del registro \texttt{CMCON}, estas terminales operan como lecturas digitales discretas provenientes de un arreglo de conmutación.

La distribución estructural del firmware se divide en:
\begin{enumerate}
    \item \textbf{Estado de Inactividad (Reposo):} Si no se detecta ningún estímulo alto en las entradas, el Puerto B se fuerza a un estado nulo (\texttt{0x00}), extinguiendo la iluminación y pausando el flujo mediante una subrutina de retención anti-rebote.
    \item \textbf{Efecto Estroboscópico (Secuencia 1):} Conmuta rítmicamente bloques de 4 bits (nibbles) con temporizaciones de $300\text{ ms}$, alternando de forma simétrica entre \texttt{0xF0} y \texttt{0x0F}.
    \item \textbf{Barrido Lateral Inyectado (Secuencia 2):} Desplaza un bit encendido de derecha a izquierda saltando de forma deliberada el canal \texttt{RB4} mediante la asignación directa del valor \texttt{0x20} posterior a \texttt{0x08}.
    \item \textbf{Barrido de Fase Alta (Secuencia 3):} Variante indexada que arranca la propagación del bit en la zona superior de la barra (\texttt{RB5}) y desciende de manera cíclica hacia los primeros canales.
\end{enumerate}

\subsection{Código de Control en C}
\begin{lstlisting}[language=C, caption={Firmware del Secuenciador para la barra de LEDs}, label={code:practica8}]
#pragma config FOSC = XT       // Oscilador de cristal externo estandar
#pragma config WDTE = OFF      // Desactivacion del Watchdog Timer
#pragma config PWRTE = ON      // Habilitacion del temporizador de encendido
#pragma config MCLRE = ON      // Master Clear habilitado externamente
#pragma config BOREN = ON      // Reset por caida de tension activo
#pragma config LVP = OFF       // Programacion en bajo voltaje inhabilitada
#pragma config CPD = OFF       // Proteccion de datos EEPROM nula
#pragma config CP = OFF        // Proteccion de codigo de programa nula

#define _XTAL_FREQ 4000000     // Frecuencia operativa de 4MHz para macros de retraso

#include <xc.h>

// Definicion de modulos de software
void suspender_flujo(void);
void configurar_perifericos(void);
void ejecutar_patron_a(void);
void ejecutar_patron_b(void);
void ejecutar_patron_c(void);

void main(void) {
    configurar_perifericos();
    
    while(1) {
        // Validacion de estado neutro: sin entradas activas
        if(RA0 == 0 && RA1 == 0 && RA2 == 0) {
            PORTB = 0x00;
            suspender_flujo(); 
        }
        
        // Disparo selectivo de rutinas
        if(RA0 == 1 && RA1 == 0 && RA2 == 0) { ejecutar_patron_a(); }
        if(RA0 == 0 && RA1 == 1 && RA2 == 0) { ejecutar_patron_b(); }
        if(RA0 == 0 && RA1 == 0 && RA2 == 1) { ejecutar_patron_c(); }
    }
}

void suspender_flujo(void) {
    while(RA0 == 0 && RA1 == 0 && RA2 == 0) {
        // Bucle de espera pasiva hasta cambio de estado
    }
    __delay_ms(100); // Mitigacion de transitorios mecanicos
}

void configurar_perifericos(void) {
    CMCON  = 0x07;  // Inactivacion de comparadores, pines en modo digital
    TRISA  = 0x07;  // Terminales RA0, RA1, RA2 configuradas como entradas
    TRISB  = 0x00;  // Puerto B completo actuando como bus de salida
    PORTA  = 0x08;  // Estado inicial preventivo para RA3
    PORTB  = 0x00;  // Inicializacion de salidas en estado bajo
}

void ejecutar_patron_a(void) {
    PORTB = 0xF0; __delay_ms(300);
    PORTB = 0x0F; __delay_ms(300);
}

void ejecutar_patron_b(void) {
    PORTB = 0x01; __delay_ms(300);
    PORTB = 0x02; __delay_ms(300);
    PORTB = 0x04; __delay_ms(300);
    PORTB = 0x08; __delay_ms(300);
    PORTB = 0x20; __delay_ms(300); // Bifurcacion: Omision selectiva de RB4
    PORTB = 0x40; __delay_ms(300);
    PORTB = 0x80; __delay_ms(300);
}

void ejecutar_patron_c(void) {
    PORTB = 0x20; __delay_ms(300);
    PORTB = 0x40; __delay_ms(300);
    PORTB = 0x80; __delay_ms(300);
    PORTB = 0x01; __delay_ms(300);
    PORTB = 0x02; __delay_ms(300);
    PORTB = 0x04; __delay_ms(300);
    PORTB = 0x08; __delay_ms(300);
}
\end{lstlisting}

\subsection{Conclusiones}
La implementación física ratificó la viabilidad del microcontrolador PIC16F877A para conmutar flujos de potencia digitales en respuesta a interacciones externas[cite: 37]. Se concluye que la correcta manipulación de registros estructurales como \texttt{TRIS} y \texttt{CMCON} es imperativa; omitir la desconfiguración analógica provocaría lecturas flotantes erráticas en el bus de entrada[cite: 38, 39]. Asimismo, la modularización del firmware a través de funciones dedicadas optimiza la legibilidad del código y mitiga de forma eficiente el impacto de los rebotes mecánicos en las terminales activas[cite: 40, 41].
